% Irodov Part III §3.5 problems 263--274, PDF pp.139--140,
% printed pp.130--131; figures 3.73--3.74 inspected.
\electromagproblem{263}{130}
很长直螺线管单位长度匝数为 $n$，
电流为 $I$。求绕组侧面受到的磁压力。
\exsolution
长螺线管内外磁场分别近似为
$B_{\rm in}=\mu_0nI$ 与 $B_{\rm out}=0$。
侧面磁压力等于两侧磁能密度差：
\[
p_m=\frac{B_{\rm in}^2-B_{\rm out}^2}{2\mu_0}.
\]
故
\[
\boxed{p_m=\frac{\mu_0n^2I^2}{2}}.
\]
\excheck 该压力沿径向\emph{向外}，
螺线管绕组受膨胀张力，
不同于通轴向电流的薄壁圆柱箍缩。

\electromagproblem{264}{130}
长单层螺线管半径 $R$，
单位长度匝数 $n$，
通电流 $I$。
若每一匝导线能够承受的最大拉力为
$F_{\rm lim}$，求断裂临界电流。
\exsolution
单匝在自身之外由其他匝产生的有效磁场
近似是长管内部场的一半：
$B_{\rm eff}\simeq\mu_0nI/2$。
一匝导线单位长度的径向力为
$f=IB_{\rm eff}=\mu_0nI^2/2$。
圆环半径 $R$，用一小段弧的径向受力平衡：
$T\,\mathrm d\theta=fR\,\mathrm d\theta$，
所以
$T=\mu_0nRI^2/2$。
当 $T=F_{\rm lim}$ 时，
\[
\boxed{I_{\rm lim}=
\sqrt{\frac{2F_{\rm lim}}{\mu_0nR}}}.
\]
\excheck $I_{\rm lim}$ 必须带\emph{平方根}，
否则 $2F_{\rm lim}/(\mu_0nR)$
的单位是安培的平方而不是安培。
原书答案纯文本抽取有漏根号风险。

\electromagproblem{265}{130}
平板电极面积 $S$、间距 $d$，
浸没在电阻率为 $\rho$、
以速率 $v$ 平行极板流动的液态导体中。
外磁场 $B$ 在板平面内垂直于流速，
电极以电阻 $R$ 闭合。
求电阻负载上的功率及其最大值。
\exsolution
液体中的动生电场
$\bm E_{\rm mot}=\bm v\times\bm B$，
其跨板距离产生电动势
$\mathcal E_{\rm mot}=vBd$。
液体内部沿板法线的电阻为
$R_{\rm in}=\rho d/S$。
因此负载电流
$I=\mathcal E_{\rm mot}/(R+R_{\rm in})$，
负载功率
\[
\boxed{P_R=
\frac{v^2B^2d^2 R}{(R+\rho d/S)^2}}.
\]
最大功率传输条件
$R=R_{\rm in}$，
于是
\[
\boxed{R_*=\frac{\rho d}{S}},
\quad
\boxed{P_{\max}=
\frac{v^2B^2dS}{4\rho}}.
\]
\excheck 这是液体磁流体发电的
简化模型，忽略涡流与边缘效应；
最大功率工况下内外电阻相等。

\electromagproblem{266}{130}
圆截面铜导线半径
$R=5.0\,\mathrm{mm}$，
沿轴通电流 $I=50\,\mathrm A$。
铜中载流电子数密度
$n=0.9\times10^{23}\,\mathrm{cm^{-3}}$。
求轴线与外表面之间因自磁场产生的横向电势差。
\exsolution
铜线中电子漂移速度模
$v_d=I/(n|e|\pi R^2)$。
导线内部自磁场随半径变化：
$B_\varphi(r)=\mu_0Ir/(2\pi R^2)$。
为平衡漂移电子的径向磁力，
霍尔电场大小
$|E_r|=v_d B_\varphi(r)$。
积分：
\[
|\Delta V|=\int_0^R|E_r(r)|\,\mathrm dr
=\frac{\mu_0Iv_d}{4\pi}.
\]
代漂移速度：
\[
\boxed{|\Delta V|=
\frac{\mu_0I^2}{4\pi^2R^2n|e|}}.
\]
按给定数据，$n=9\times10^{28}\,\mathrm{m^{-3}}$，
数值
\[
\boxed{|\Delta V|\simeq2.2\times10^{-10}\,\mathrm V
=220\,\mathrm{pV}}.
\]
\excheck 本版书末简答数值似为
“$2\,\mathrm{pV}$”，
但其写出的同一解析公式与原题数值
直接代入得到约220pV，
相差约两个数量级；待核查原版物性数值
或单位排印。

\electromagproblem{267}{130--131}
钠导体霍尔测量中，
横向电场
$E_H=5.0\,\mathrm{\mu V/cm}$，
纵向电流密度
$j=200\,\mathrm{A/cm^2}$，
磁感应强度 $B=1.00\,\mathrm T$。
求自由电子数密度及每原子对应自由电子数。
\exsolution
单载流子霍尔平衡
$|e|E_H=|e|v_dB$，
又 $j=n|e|v_d$，
所以
\[
\boxed{n=\frac{jB}{|e|E_H}}.
\]
将 $j=2.0\times10^6\,
\mathrm{A/m^2}$、
$E_H=5.0\times10^{-4}\,\mathrm{V/m}$，
得
\[
\boxed{n\simeq2.5\times10^{28}\,\mathrm{m^{-3}}}.
\]
钠金属原子数密度估计为
$n_{\rm atom}=\rho_{\rm Na}N_A/M_{\rm Na}$，
取 $\rho_{\rm Na}\simeq970\,\mathrm{kg/m^3}$、
$M_{\rm Na}\simeq0.023\,\mathrm{kg/mol}$，
得
$n_{\rm atom}\simeq2.5\times10^{28}\,\mathrm{m^{-3}}$。
因此
$\boxed{n/n_{\rm atom}\simeq1.0}$。
\excheck 原书书末纯文本提取
可能将数密度单位误作 $\mathrm{m^{-1}}$，
正确单位必须是 $\mathrm{m^{-3}}$。

\electromagproblem{268}{131}
铜导体在 $B=100\,\mathrm{mT}$ 下
霍尔测量，横向霍尔电场比纵向电场
小 $\eta=3.1\times10^3$ 倍。
求导电电子的漂移迁移率 $\mu$。
\exsolution
单载流子近似
$E_H=v_dB$，
纵向漂移规律
$v_d=\mu E_\parallel$。
故两场比
$E_H/E_\parallel=\mu B=1/\eta$。
所以
\[
\boxed{\mu=\frac1{\eta B}}.
\]
代 $B=0.100\,\mathrm T$、
$\eta=3.1\times10^3$，
\[
\boxed{\mu\simeq3.2\times10^{-3}\,
\mathrm{m^2/(V\,s)}}.
\]
\excheck \textbf{原书书末数值冲突：}
该版简答似写
$3.2\times10^{-2}\,\mathrm{m^2/(V\,s)}$，
但同一公式 $1/(\eta B)$ 代入0.1T
得到 $3.2\times10^{-3}$，
相差十倍，已保留计算过程。

\electromagproblem{269}{131}
长直导线通电流 $I$，
距它 $r$ 处有磁矩为 $\bm p_m$
的小电流环。
求电流环所受力的模与方向，
当磁矩：(a) 平行直线导体；
(b) 沿径向 $\bm e_r$；
(c) 与当地磁场方向相同。
\exsolution
长线磁场
$\bm B(r)=k_m r^{-1}\bm e_\varphi$，
$k_m=\mu_0I/(2\pi)$。
固定磁矩的力可用
$\bm F=(\bm p_m\cdot\nabla)\bm B$。
(a) $\bm p_m=p_m\bm e_z$，
场与轴向坐标无关，
$\boxed{\bm F_a=0}$。
(b) $\bm p_m=p_m\bm e_r$，
径向导数
$\partial_r\bm B=-k_mr^{-2}\bm e_\varphi$，
\[
\boxed{\bm F_b=
-\frac{\mu_0Ip_m}{2\pi r^2}\bm e_\varphi}.
\]
(c) $\bm p_m=p_m\bm e_\varphi$，
沿方位角方向移动时
$\partial_\varphi\bm e_\varphi=-\bm e_r$，
定向导数
$(p_m/r)\partial_\varphi\bm B
=-p_mk_m r^{-2}\bm e_r$，
故
\[
\boxed{\bm F_c=
-\frac{\mu_0Ip_m}{2\pi r^2}\bm e_r}.
\]
\excheck (b) 的力与当地 $\bm B$ 反向；
(c) 指向导线轴；
符号随磁矩或电流反向而改变。

\electromagproblem{270}{131}
半径 $R$、电流 $I$ 的圆环轴线上
距中心 $x$ 处放一只小线圈，
磁矩 $p_m$ 平行圆环轴。
求线圈所受磁力模。
\exsolution
大圆环轴向场
$B(x)=\mu_0IR^2/
[2(R^2+x^2)^{3/2}]$。
小磁偶极子在非均匀场中受力
$F_x=p_m\,\mathrm dB/\mathrm dx$，
其模
\[
|F_x|=p_m\frac{3\mu_0IR^2|x|}
{2(R^2+x^2)^{5/2}}.
\]
故
\[
\boxed{F=
\frac{3\mu_0IR^2p_m|x|}
{2(R^2+x^2)^{5/2}}}.
\]
\excheck $x=0$ 因对称性受力为零，
远场 $x\gg R$ 时按 $x^{-4}$ 衰减。

\electromagproblem{271}{131}
两个很小的同轴线圈磁矩为
$p_{m1}=4.0\,\mathrm{mA\,m^2}$、
$p_{m2}=6.0\,\mathrm{mA\,m^2}$，
间距 $\ell=20\,\mathrm{cm}$
远大于线圈大小。
求相互作用力大小。
\exsolution
磁偶极子1的轴上远场
$B(\ell)=\mu_0p_{m1}/(2\pi\ell^3)$。
第二个同轴磁偶极子的力模
\[
F=p_{m2}\left|\frac{\mathrm dB}{\mathrm d\ell}\right|
=\frac{3\mu_0p_{m1}p_{m2}}{2\pi\ell^4}.
\]
代数值：
\[
\boxed{F\simeq9.0\times10^{-9}\,\mathrm N=9\,\mathrm{nN}}.
\]
\excheck 两磁矩同向时吸引，
反向时排斥；题目未给相对方向时
所列为力的\emph{大小}。

\electromagproblem{272}{131}
薄圆盘形永磁体沿轴磁化，
盘半径 $R=1.0\,\mathrm{cm}$。
在轴线上距离 $x=10\,\mathrm{cm}$
处测得磁场 $B=30\,\mathrm{\mu T}$。
估算沿圆盘边缘流动的等效分子电流 $I'$。
\exsolution
薄轴向磁化圆盘的边缘束缚电流
可等效为半径 $R$ 的电流环，
轴上磁场
$B(x)=\mu_0I'R^2/
[2(R^2+x^2)^{3/2}]$。
当 $x\gg R$，
$B\simeq\mu_0I'R^2/(2x^3)$，
故
\[
\boxed{I'\simeq\frac{2Bx^3}{\mu_0R^2}}.
\]
代入 SI 数值：
\[
\boxed{I'\simeq4.8\times10^2\,\mathrm A
\approx0.5\,\mathrm{kA}}.
\]
\excheck 这是\emph{磁化分子束缚电流}的
宏观等效，不是外电源实际输送的自由电流。

\electromagproblem{273}{131}
均匀各向同性磁介质和真空
之间有平界面，介质相对磁导率 $\mu$。
真空侧磁感应强度模为 $B$，
与界面法线夹角 $\alpha$。
求介质内近界面场强模 $B'$。
\exsolution
若界面无自由面电流，
磁感应强度的法向分量连续：
$B'_n=B_n=B\cos\alpha$。
磁场强度的切向分量连续：
$H'_t=H_t=B\sin\alpha/\mu_0$。
因此介质内
$B'_t=\mu\mu_0H'_t=\mu B\sin\alpha$。
模的平方为
$B'^2=B_n'^2+B_t'^2$，
所以
\[
\boxed{B'=B\sqrt{\cos^2\alpha+
\mu^2\sin^2\alpha}}.
\]
\excheck $\alpha=0$ 时 $B'=B$，
$\alpha=\pi/2$ 时 $B'=\mu B$，
分别是法向与切向边界条件的直接极限。

\electromagproblem{274}{131--132}
图3.74：磁介质相对磁导率 $\mu$，
与真空以平面分界。
真空侧磁感应强度大小 $B$，
与法线 $\bm n$ 夹角 $\theta$。
(a) 求以界面上一点为球心、半径 $R$
的球面上 $\bm H$ 的通量；
(b) 求图示横跨界面的边长 $\ell$
的正方形闭路上 $\bm B$ 的环流。
\exsolution
法向 $B_n=B\cos\theta$ 连续，
所以
$H_{{\rm vac},n}=B\cos\theta/\mu_0$，
$H_{{\rm mag},n}=B\cos\theta/(\mu\mu_0)$。
(a) 两半球的法向投影面积均为
$\pi R^2$，法线对整个闭合球面
方向相反，净通量
\[
\boxed{\oint_S\bm H\cdot\mathrm d\bm S
=\frac{\pi R^2B\cos\theta}{\mu_0}
\left(1-\frac1\mu\right)}.
\]
(b) 磁场切向分量在两侧为
$B_{{\rm vac},t}=B\sin\theta$、
$B_{{\rm mag},t}=\mu B\sin\theta$。
沿图示正向闭合方框，平行边贡献相减，
垂直边因边长趋零或场分量相消，
得到
\[
\boxed{\oint_\Gamma\bm B\cdot\mathrm d\bm\ell
=(1-\mu)B\ell\sin\theta}.
\]
\excheck 不能把 $\bm B$ 的环流
与 $\bm H$ 的环流混淆：
前者可因束缚电流而非零。
